\section{Problem}\label{sec:background}

\paragraph*{MDPs and requirements}
An instance is an MDP $M = (S, s_0, A, P)$ given as a PRISM model, with \emph{policy actions} $A$ (action labels the
controller chooses) and requirements $\varphi_i \bowtie_i \theta_i$: a PRISM path formula $\varphi_i$ (reachability,
until, or an LTL formula such as ``goals 1 and 2 in order before goal 3 is ever entered''), a bound
${\bowtie_i} \in \{\geq, \leq\}$ and a threshold $\theta_i$. A \emph{family} is a set of such instances that share
the policy-visible variables $X \subseteq$ state variables, the actions and the requirement names; their state spaces,
constants and thresholds may differ (a $9\times 9$ and a $20\times 20$ grid are members of one family).

\paragraph*{Rule lists}
A controller is an ordered list $R = \langle (c_1, a_1), \dots, (c_n, a_n) \rangle$ of rules, where $c_j$ is a
Boolean condition over $X$, the instance's named constants $K$ and the domain's features (conditions written in the
same language), and $a_j \in A \cup \{\any\}$. Per-action features (``\texttt{wall}'', ``\texttt{closer}'') are
evaluated for the rule's action; a rule $(c, \any)$ allows every action $b$ for which $c[b]$ holds. In a state $s$,
the first rule that allows some action decides, and its allowed actions are the controller's choices; if no rule
decides, every action is allowed. Conditions may only use the numbers 0 and 1, so a rule names instance values
through $K$ and means something on every member (\emph{general mode}).

\paragraph*{Verification and certification}
Restricting $M$ to the allowed actions gives the induced MDP $M_R$. Its best and worst case for a requirement are the
optimum over all schedulers in either direction; PRISM's schedulers may use memory, randomisation and the full state,
so the interval $[w_i, b_i]$ contains the value of \emph{every} deployment consistent with $R$. If every $w_i$ meets
its bound, every such deployment meets every requirement: $R$ is \emph{certified} on $M$. Certification is per
instance and exact (we use interval iteration for the final verdicts); a rule list certified on 30 instances says
nothing about a 31st, and we never claim it does.

\paragraph*{Task}
Given a few small training instances of a family, write one rule list $R$, then freeze it and certify it on as many
held-out instances as possible, including larger ones. Unlike optimisation, a held-out instance counts only if every
requirement holds in the worst case.
